\scanpage{131}
\mspage{355}

the set of fixed \unsure{points} of $U_{0,3}$ under $\dot{\sigma}_\infty$ is reduced to a single element, which also implies that
\deleted{\illegible{1}}
\[
\Pi_{0,3}^{\sigma_\infty}=\{1\}.
\]

We shall \unsure{also} assume that these facts \unsure{extend} to the elements of $\widehat{\Pi}_{0,3}$ --- which will be the case if (in addition to what has been assumed about it with regard to $\rho$) we assume that we \unsure{still} have
\[
\tag{38}
\widehat{\Pi}_{0,3}.
\]

Note that here we have
\[
\tag{39}
G^{\sigma_\infty}=\mathcal T_{0,3}^{\sigma_\infty}\simeq\mathbb Z/2\mathbb Z\times\mathbb Z/2\mathbb Z,
\]
subgroup generated by $\dot{\sigma}_\infty,\unsure{\dot\tau|\tau}$, whose non\unsure{zero} elements are $\dot{\sigma}_\infty,\tau,\dot{\tau}'_\infty$ --- which \unsure{lift} to $\sigma_\infty,\tau_\infty,\tau'_\infty$ ; and we have
\[
\tag{40}
\mathcal E^{\sigma_\infty}=\mathcal G_{0,3}^{\sigma_\infty}\simeq\mathbb Z/2\mathbb Z\times\mathbb Z/2\mathbb Z,
\]
group generated by $\sigma_\infty,\tau_\infty$ (the other nontrivial element being $\tau'_\infty$). If then we write a solution $h\inserted{\in\pi}$ (of (35) in the form (36)), with $\beta\in\mathcal G_{0,3}$ only (not only $\beta\in\pi$), then the indeterminacy is in the group (40) of order (4) ; More precisely, if $\beta_0$ is the solution in $\pi$, then the other solutions are of the form
\[
\tag{41}
\beta=\beta_0\sigma_\infty^\varepsilon\tau_\infty^{\varepsilon'}
\qquad \varepsilon,\varepsilon'\in\mathbb Z/2\mathbb Z
\]

\margin{130}

\scanpage{132}

(with $\varepsilon,\varepsilon'$ uniquely determined). In order to have $\beta\in\mathcal G_{0,3}^{+}$, it is necessary and sufficient that $\varepsilon'=0$ !

\margin{131}

\scanpage{133}
\mspage{356}

\margin{These notations will be \unsure{revised} once again in \S\ 46}

\uline{44}\qquad \uline{Return to the notation.} \uline{Setting up equations} \uline{in} $\mathrm{Sl}(2,\widehat{\mathbb Z})$.

Let $u$ be an automorphism of $\widehat{\Pi}_{0,3}$ \inserted{commuting up to inner automorphisms : $\unsure{\mathcal G_3}$\deleted{\illegible{1}}} --- or what amounts to the same, an automorphism of $\widehat{\mathcal G}_{0,3}^{+}$ which induces the identity in $\mathcal T_{0,3}^{+}=\unsure{\mathcal G_3}$. Then $u$ is known when we know its actions on the two generators $\sigma_\infty,\rho$ of $\widehat{\mathcal G}_{0,3}^{+}$, which will resp. be elements of $\widehat{\mathcal G}_{0,3}^{+}$ over $\dot{\sigma}_\infty$ and $\dot\rho$,

namely
\[
\tag{1}
\boxed{u(\sigma_\infty)=h\sigma_\infty,\qquad u(\rho)=g\rho}
\qquad (h,g\in\widehat{\Pi}_{0,3})
\]
where $h,g$ \uline{are subject to the conditions}
\[
\tag{2}
\boxed{h\sigma_\infty(h)=1\qquad g\rho(g)\rho^2(g)=1},
\]
expressing respectively the relations
\[
\tag{3}
\boxed{(h\sigma_\infty)^2=(g\rho)^3=1}.
\]

(There is in addition the \unsure{preliminary} condition to check --- namely that (3) are defining relations for $\widehat{\mathcal G}_{0,3}^{+}$ \dots).

\deleted{\illegible{1}}\inserted{For the moment}, there is no need to impose any relation between $g$ and $h$ --- except the relations (2) on each of these elements --- since $\widehat{\mathcal G}_{\unsure{0,3}}^{+}$ is the group with defining relations
\[
\sigma_\infty^2=\rho^3=1.
\]

\margin{132}

\scanpage{134}

The elements $g$ and $h$, thus taken \unsure{as} automorphisms of $\widehat{\Pi}_{0,3}$, express the commutation relations of\deleted{s} $u$ with \deleted{$\sigma_\infty$} $\rho$ and $\sigma_\infty$ respectively :
\[
\tag{4}
[u,\rho]=\underbrace{u\rho u^{-1}}_{u(\rho)}\rho^{-1}=\operatorname{int}(g),
\qquad
[u,\sigma_\infty]=u\sigma_\infty u^{-1}\sigma_\infty=\operatorname{int}(h)
\]
--- it is a remarkable thing that knowledge of the elements $g,h$ of $\pi$ satisfying (4), already \uline{\unsure{implies|imply}} knowledge of $u$ (\unsure{this|which} \unsure{to the} fact that $\sigma_\infty,\rho$ generate $\mathcal G_{0,3}^{+}$ \dots).

Let us now write the compatibility of $u$ with the loop structure, \deleted{\illegible{3}}\inserted{We had it} \deleted{\unsure{formulated}}\inserted{\unsure{expressed|expressed}} previously by a relation of the form
\[
u(\varepsilon_0)=\operatorname{int}(g_0)(\varepsilon_0^p)
\quad\text{or}\quad
u(\ell_0)=\operatorname{int}(g_0)\ell_0^p,
\]
for
\[
p\in\widehat{\mathbb Z}^{*},\qquad g_0\in\widehat{\Pi}_{0,3}
\]
suitably chosen. The notation $g_0$ is not all that happy, since $g$ and $g_0$ play quite dissimilar roles. The subscript $0$ in $g_0$ was intended to allow us to write, more generally
\[
\deleted{\unsure{g_i}}\quad
u(\ell_i)=\operatorname{int}(g_i)\ell_i^p
\qquad i\in\{0,1,\infty\}.
\]

\margin{133}

\scanpage{135}
\mspage{357}

I prefer to use the letter $f_0$ --- with the option of dropping the subscript if no confusion is possible. The letter $p$ for a multiplier is not \inserted{a choice} that is very happy, for later on \deleted{\unsure{$p$}} we shall \inserted{\unsure{surely}} need this letter to denote a characteristic --- I prefer the letter $\mu$ (Greek initial of ``multiplier''), and the loop relation \inserted{\unsure{would}} therefore take the form
\[
\deleted{\unsure{(5)}}\qquad
u(\varepsilon_0)=\operatorname{int}(f_0)\varepsilon_0^\mu.
\qquad
\deleted{\unsure{\text{or}}\quad u(\ell_0)=\operatorname{int}(f_0)}
\]

In fact, I prefer to replace $f_0$ in it by $f_0^{-1}$, to obtain later on formulas in which the symmetry of the roles played by $\rho$ and by $\sigma_\infty$ \unsure{appears} more clearly. --- in particular I would like the formula defining $h_0$ in terms of $g_0$ (\S\ 43, formula (65,)
\[
h_0=g_0^{-1}\sigma_\infty(g_0)
\]
to become instead
\[
h_0=g_0\sigma_\infty(g_0)^{-1},
\]
in harmony with
\[
\xi=\lambda\rho(\lambda)^{-1}\quad\dots
\]

Thus the ``\unsure{loop|loop} relation'' takes the form
\[
\tag{5}
u(\varepsilon_0)=\operatorname{int}(f_0^{-1})\varepsilon_0^\mu
\qquad f_0\in\widehat{\Pi}_{0,3},\ \mu\in\widehat{\mathbb Z}^{*},
\]
\deleted{\unsure{rather} \unsure{we} immediately deduce}
\[
\tag{6}
u(\ell_0)=\operatorname{int}(f_0^{-1})\ell_0^\mu,
\]
from which we immediately deduce

\margin{134}

\scanpage{136}

\[
\tag{7}
\left\{
\begin{aligned}
u(\varepsilon_1)&=\operatorname{int}(f_1^{-1})(\varepsilon_1^\mu),
&u(\ell_1)&=\operatorname{int}(f_1^{-1})(\ell_1^\mu),\\
u(\varepsilon_\infty)&=\operatorname{int}(f_\infty^{-1})(\varepsilon_\infty^\mu),
&u(\ell_\infty)&=\operatorname{int}(f_\infty^{-1})(\ell_\infty^\mu),
\end{aligned}
\right.
\]
\unsure{where|or}
\[
\tag{8}
f_1=\rho(f_0)g^{-1},\qquad
f_\infty=\rho(f_1)g^{-1}=\rho^2(f_0)\rho(g)^{-1}g^{-1}.
\]

In what follows, we shall simply write $f$ instead of $f_0$, and we are going to transform equation (5) \deleted{\illegible{1}} into
\[
\tag{8bis}
\operatorname{int}(f)u(\varepsilon_0)=\varepsilon_0^\mu,
\]
by making the left-hand side explicit
\[
\begin{aligned}
\operatorname{int}(f)u(\varepsilon_0)
&=f\,u(\sigma_\infty\rho)f^{-1}
=f\underbrace{u(\sigma_\infty)}_{h\sigma_\infty}\underbrace{u(\rho)}_{g\rho}f^{-1}
=fh\sigma_\infty g\rho f^{-1}\\
&=fh\sigma_\infty g\sigma_\infty^{-1}\underbrace{\sigma_\infty\rho}_{\varepsilon_0}f^{-1}
=fh\sigma_\infty(g)\underbrace{\varepsilon_0f^{-1}\varepsilon_0^{-1}}\varepsilon_0
\end{aligned}
\]
i.e.
\[
\tag{9}
\operatorname{int}(f)u(\varepsilon_0)
=\bigl[fh\sigma_\infty(g)(\sigma_\infty\rho)(f)^{-1}\bigr]\varepsilon_0
=\bigl[f,h\sigma_\infty(g\rho(f)^{-1}\bigr]\varepsilon_0
\]
where the \unsure{four|three} factors $f,h,\sigma_\infty(g\rho(f)^{-1}$ in $[\ ]$ are $\in\widehat{\Pi}_{0,3}$. Thus \unsure{(5)|(9)} or (8), takes the form
\[
fh\sigma_\infty(g\rho(f)^{-1})=\varepsilon_0^{\mu-1}\,\deleted{\illegible{1}}
\]
which gives, setting
\[
\tag{10}
\mu-1=2\nu\qquad(\nu\in\widehat{\mathbb Z})
\]
(so we write $\nu$ for what was denoted by $\gamma$ previously) we find the relation
\[
fh\sigma_\infty\bigl(g\rho(f)^{-1}\bigr)=\ell_0^\nu
\]
or again, transforming by $\sigma_\infty$
\[
\sigma_\infty(\unsure{f})\sigma_\infty(h)g\rho(f)^{-1}=\ell_1^\nu
\]
; that is

\margin{135}

\scanpage{137}
\mspage{358}

\[
\boxed{\sigma_\infty(h).g=\sigma_\infty(f)^{-1}\ell_1^\nu\rho(f)}.
\qquad (11)
\]

This relation, which relates $h$ and $g$ via $f$, can rightly be called the ``loop relation''.
We see, from this relation, that \deleted{\illegible{1}} not only does knowledge \inserted{of $u$ i.e.} \unsure{of the|of} $h$ and $g$ determine $f$ and $\mu$ i.e. $f$ and $\nu$, but also conversely, once $f$ and $\nu$ are known, $g$ and $h$ determine each other via (11).

\margin{$f$ determined modulo $L_0$ only $f\mapsto\ell_0^n f$ $(n\in\widehat{\mathbb Z})$.}

Note that for $u$ corresponding to : \deleted{\unsure{$f$|$\xi$}}, $g,h,f$, if $\gamma\in\widehat{\Pi}_{0,3}$, the \unsure{automorphism} \inserted{$u'=$} $\operatorname{int}(\gamma)\circ u$ corresponds to
\[
g'=\gamma g\rho(\gamma^{-1}),\quad
h'=\gamma h\sigma_\infty(\gamma)^{-1},\quad
f'=f\gamma^{-1}.
\qquad (12)
\]

\margin{This holds more \unsure{generally} if $\gamma\in\unsure{\widetilde{\mathcal G}_{0,3}|\mathcal G_{0,3}}$, and even for $\gamma\in\unsure{\widehat{\widetilde{\mathcal G}}_{0,3}|\widehat{\mathcal G}_{0,3}}$.}

The equations (2) ``in $h$'' and ``in $g$'' are solved parametrically in a unique way by
\[
g=\alpha\rho(\alpha)^{-1},\quad
h=\beta\sigma_\infty(\beta)^{-1}
\qquad (13)
\]
where
\[
\left\{\begin{array}{l}
\alpha\in\widehat{\widetilde{\Pi}}_{0,3}=\widehat{\Pi}_{0,3}.\{1,\tau_\infty\}\\
\beta\in\widehat{\Pi}_{0,3}
\end{array}\right.
\qquad (14)
\]
to be precise, giving $g,h$ satisfying (2) is equivalent resp. to giving $\alpha,\beta$ in $\widehat{\widetilde{\Pi}}_{0,3}$ and $\widehat{\Pi}_{0,3}$, via the formulas (13).

\margin{136}

\scanpage{138}

\inserted{The elements} $g$ and $h$ being thus made explicit in terms of $\alpha$ and $\beta$, the \inserted{loop} relation (11) can be considered as a relation between $\alpha,\beta,f,\nu$
\[
\sigma_\infty(\beta\sigma_\infty(\beta)^{-1})(\alpha\rho(\alpha)^{-1})
=\sigma_\infty(f)^{-1}\ell_1^\nu\rho(f),
\qquad (15)
\]
where \deleted{$\nu,f$ \unsure{are}} $\nu$ is uniquely determined by $\alpha,\beta$, and $f$ is determined in $L_0\backslash\widehat{\Pi}_{0,3}$.
(We shall note that the right-hand side of (15) does not change, on changing $f$ to $\ell_0^n f$, with $n\in\widehat{\mathbb Z}$).

The main asymmetry, it seems to me, between $\sigma_\infty$ and $\rho$ in this theory, is that \deleted{\illegible{1}} we always have $\beta\in\widehat{\Pi}_{0,3}$, whereas $\alpha$ is not necessarily in $\widehat{\Pi}_{0,3}$ -- it is, \unsure{as we have seen}, iff $\nu\equiv0\ (3)$ i.e. $\mu\equiv1\ (3)$.

\deleted{So if we use \illegible{1} $\operatorname{int}(\beta^{-1})$ \illegible{1}}
The automorphism
\[
u'=\operatorname{int}(\beta^{-1})u
\]
will therefore be \unsure{congruent} to $u$ modulo \unsure{inner} automorphisms, and will correspond to
\[
\left\{\begin{array}{l}
g'=\beta^{-1}g\rho(\beta)=(\beta^{-1}\alpha)\rho(\beta^{-1}\alpha)^{-1}\\
h'=\beta^{-1}h\sigma_\infty(\beta)=1\\
f'=f\beta
\end{array}\right.
\qquad (16)
\]
-- thus $u$ is \deleted{conjugate} \inserted{congruent modulo} inner automorphisms

\margin{137}